\documentclass[11pt,a4paper]{ctexart}
\usepackage[margin=2.2cm]{geometry}
\usepackage{amsmath,amssymb}
\usepackage{graphicx}
\usepackage{booktabs}
\usepackage{enumitem}
\usepackage{xcolor}
\usepackage{hyperref}
\usepackage{array}
\usepackage{tcolorbox}
\tcbuselibrary{skins,breakable}

\hypersetup{
  colorlinks=true,
  linkcolor=blue!55!black,
  urlcolor=blue!55!black
}

\setlist[itemize]{leftmargin=2em,itemsep=2pt,topsep=3pt}
\setlist[enumerate]{leftmargin=2.2em,itemsep=3pt,topsep=3pt}
\renewcommand{\arraystretch}{1.25}

\definecolor{mainblue}{RGB}{26,82,118}
\definecolor{softblue}{RGB}{235,245,252}
\definecolor{warn}{RGB}{163,72,58}
\definecolor{softwarn}{RGB}{252,240,237}
\definecolor{softgreen}{RGB}{238,248,241}

\newtcolorbox{focusbox}[1]{
  enhanced,
  breakable,
  colback=softblue,
  colframe=mainblue,
  boxrule=0.8pt,
  arc=2mm,
  left=8pt,
  right=8pt,
  top=6pt,
  bottom=6pt,
  title=\textbf{#1}
}

\newtcolorbox{warnbox}[1]{
  enhanced,
  breakable,
  colback=softwarn,
  colframe=warn,
  boxrule=0.8pt,
  arc=2mm,
  left=8pt,
  right=8pt,
  top=6pt,
  bottom=6pt,
  title=\textbf{#1}
}

\newtcolorbox{greenbox}[1]{
  enhanced,
  breakable,
  colback=softgreen,
  colframe=green!45!black,
  boxrule=0.8pt,
  arc=2mm,
  left=8pt,
  right=8pt,
  top=6pt,
  bottom=6pt,
  title=\textbf{#1}
}

\title{\heiti 第5章 理想流体力学特性\\今日复习路线}
\author{根据何枫课件整理}
\date{}

\begin{document}
\maketitle

\section*{一、今天目标}

今天不做泛泛复习，只解决一个目标：\textbf{看到第5章对应题目，知道用哪个方程、为什么能用、怎么写步骤。}

\begin{focusbox}{必须学会的四类题}
\begin{enumerate}
  \item 判断伯努利方程能不能用，常数沿哪里成立。
  \item 用非定常伯努利 / 柯西--拉格朗日积分做水柱运动、压力分布题。
  \item 用速度势处理无界膨胀球问题，例如书 4.12。
  \item 会写环量、开尔文定理、拉格朗日定理、海姆霍兹定理的条件和结论。
\end{enumerate}
\end{focusbox}

\subsection*{本次作业题目原图}

\begin{center}
\includegraphics[width=0.93\linewidth]{ch5_images/assignment_page.png}
\end{center}

\begin{warnbox}{看图时先只盯第5章}
这张图下半部分还有第6章题目。今天第5章只抓三件事：第1题环量，第2题书 4.12，第3题弯管/开阀非定常伯努利。
\end{warnbox}

\section*{二、最小公式表}

\begin{center}
\begin{tabular}{>{\raggedright\arraybackslash}p{3.5cm} >{\raggedright\arraybackslash}p{8.8cm}}
\toprule
\textbf{名称} & \textbf{公式 / 用法}\\
\midrule
环量 &
\[
\Gamma=\oint_C \vec V\cdot d\vec l
\]
\\
欧拉方程 &
\[
\frac{D\vec V}{Dt}=\vec f-\frac{1}{\rho}\nabla p
\]
\\
兰姆方程 &
\[
\frac{\partial \vec V}{\partial t}
+\nabla\frac{V^2}{2}
-\vec V\times\vec\Omega
=\vec f-\frac{1}{\rho}\nabla p
\]
\\
定常伯努利 &
\[
\frac{V^2}{2}+\frac{p}{\rho}+gz=C
\]
\\
非定常无旋伯努利 &
\[
\frac{\partial\varphi}{\partial t}
+\frac{V^2}{2}
+\frac{p}{\rho}
+gz=C(t)
\]
\\
速度势 &
\[
\vec V=\nabla\varphi
\]
\\
球对称径向流 &
\[
v_r=\frac{\partial\varphi}{\partial r}
\]
\\
\bottomrule
\end{tabular}
\end{center}

\begin{warnbox}{最容易混的点}
\begin{itemize}
  \item 理想流体固壁条件：\textbf{不可穿透}，但\textbf{允许滑移}。
  \item 定常伯努利一般沿流线成立；如果无旋，则可在全流场任意两点使用。
  \item 非定常无旋流必须多写一项：
  \[
  \frac{\partial\varphi}{\partial t}
  \]
  \item 跨越激波、强损失、粘性主导流动时，不能直接套理想伯努利。
\end{itemize}
\end{warnbox}

\section*{三、题型 1：环量题}

\subsection*{识别}

题目给速度场，又给一条封闭路径，问：

\[
\Gamma=\oint_C \vec V\cdot d\vec l
\]

\subsection*{标准写法}

在极坐标中：

\[
\vec V\cdot d\vec l=V_r\,dr+V_\theta r\,d\theta
\]

如果题目给：

\[
V_r=ar,\qquad V_\theta=0
\]

圆弧段上：

\[
dr=0,\qquad V_\theta=0
\]

所以圆弧段环量为零。

两条径向段：

\[
\Gamma_1=\int_{R_1}^{R_2}ar\,dr
=\frac{a}{2}(R_2^2-R_1^2)
\]

\[
\Gamma_2=\int_{R_2}^{R_1}ar\,dr
=-\frac{a}{2}(R_2^2-R_1^2)
\]

因此：

\[
\boxed{\Gamma=0}
\]

\section*{四、题型 2：伯努利适用条件}

\subsection*{定常伯努利}

若满足：

\[
\text{无粘、定常、不可压、质量力有势}
\]

则沿流线：

\[
\boxed{
\frac{V^2}{2}+\frac{p}{\rho}+gz=C
}
\]

如果进一步满足：

\[
\Omega=0
\]

即无旋，则伯努利常数可在全流场相同。

\begin{greenbox}{考试可直接写}
理想、定常、不可压、质量力有势流动中，伯努利方程沿流线成立；若流动无旋，则伯努利常数在全流场相同。
\end{greenbox}

\section*{五、题型 3：非定常水柱 / 开阀题}

\subsection*{识别}

出现这些字眼：

\begin{itemize}
  \item 不计粘性；
  \item 突然打开阀门；
  \item 水柱振荡；
  \item 求瞬时压力分布；
  \item 求自由面运动规律。
\end{itemize}

优先考虑非定常伯努利积分式。

\subsection*{课件题图：U 型管振荡}

\begin{center}
\includegraphics[width=0.93\linewidth]{ch5_images/slide_unsteady_utube.png}
\end{center}

\subsection*{非定常伯努利：表格版}

\begin{center}
\begin{tabular}{p{2.8cm} p{5.8cm} p{5.2cm}}
\toprule
\textbf{要判断的事} & \textbf{怎么写} & \textbf{作用}\\
\midrule
基本条件 &
无粘、非定常、不可压；质量力有势 \( \vec f=-\nabla \Pi \) &
决定可以从兰姆方程积分。\\
\midrule
沿流线积分 &
\[
\begin{aligned}
&\int_1^2 \frac{\partial V}{\partial t}\,dl\\
&+\left(\frac{V^2}{2}+\frac{p}{\rho}+gz\right)_2\\
&-\left(\frac{V^2}{2}+\frac{p}{\rho}+gz\right)_1=0
\end{aligned}
\]
&
用于水柱振荡、开阀瞬时压力分布。\\
\midrule
等截面管 &
\[
V_1=V_2=V(t)
\]
&
速度项常抵消。\\
\midrule
两端通大气 &
\[
p_1=p_2=p_a
\]
&
压力项常抵消。\\
\midrule
最后常见形式 &
\[
L\frac{dV}{dt}+g(z_2-z_1)=0
\]
&
把流动问题变成一维运动方程。\\
\bottomrule
\end{tabular}
\end{center}

\subsection*{U 型管振荡：结果表}

\begin{center}
\begin{tabular}{p{3cm} p{5.3cm} p{5.3cm}}
\toprule
\textbf{量} & \textbf{表达式} & \textbf{备注}\\
\midrule
坐标关系 &
\[
h=2z
\]
&
\(z\) 为一侧液面相对平衡位置位移。\\
\midrule
控制方程 &
\[
\frac{d^2z}{dt^2}+\frac{2g}{L}z=0
\]
&
简谐振动方程。\\
\midrule
初始条件 &
\[
z(0)=\frac{h_0}{2},\qquad \dot z(0)=0
\]
&
初始液面差为 \(h_0\)，初速为零。\\
\midrule
振荡规律 &
\[
z=\frac{h_0}{2}\cos\sqrt{\frac{2g}{L}}\,t
\]
&
直接背这个结果。\\
\midrule
周期 &
\[
T=2\pi\sqrt{\frac{L}{2g}}
\]
&
水柱越长，周期越大。\\
\midrule
速度 &
\[
V=\dot z
=-\frac{h_0}{2}\sqrt{\frac{2g}{L}}
\sin\sqrt{\frac{2g}{L}}\,t
\]
&
由位移对时间求导。\\
\bottomrule
\end{tabular}
\end{center}

\section*{六、题型 4：书 4.12 膨胀球}

\subsection*{课件题图：直角管开阀}

\begin{center}
\includegraphics[width=0.93\linewidth]{ch5_images/slide_right_angle_pipe.png}
\end{center}

这类题和作业第3题同源：都是等截面水柱、突然开阀、不计粘性。关键不是静水压，而是用非定常伯努利中的
\[
\int \frac{\partial V}{\partial t}\,dl
\]
处理瞬时加速度项。

\subsection*{题目模型}

原静止无界理想均质不可压流体中，有一球形物体，其半径按

\[
R=1+3t
\]

膨胀。求压力场。已知无穷远处压力为 \(p_\infty\)，密度为 \(\rho\)，不计质量力。

\subsection*{步骤 1：由不可压缩连续性求速度}

球面速度：

\[
\dot R=3
\]

球对称径向流，任意球面流量相同：

\[
4\pi r^2v_r=4\pi R^2\dot R
\]

所以：

\[
v_r=\frac{R^2\dot R}{r^2}
\]

代入 \(R=1+3t\)：

\[
\boxed{
v_r=\frac{3(1+3t)^2}{r^2}
}
\]

\subsection*{步骤 2：求速度势}

径向无旋流有：

\[
v_r=\frac{\partial\varphi}{\partial r}
\]

因此：

\[
\varphi=\int \frac{R^2\dot R}{r^2}\,dr
=-\frac{R^2\dot R}{r}
\]

代入：

\[
\boxed{
\varphi=-\frac{3(1+3t)^2}{r}
}
\]

\subsection*{步骤 3：用非定常无旋伯努利求压力}

无穷远处：

\[
v_\infty=0,\qquad \varphi_\infty=0,\qquad p=p_\infty
\]

因此：

\[
\frac{\partial\varphi}{\partial t}
+\frac{v^2}{2}
+\frac{p}{\rho}
=
\frac{p_\infty}{\rho}
\]

即：

\[
p=p_\infty-\rho
\left(
\frac{\partial\varphi}{\partial t}
+\frac{v^2}{2}
\right)
\]

计算：

\[
\frac{\partial\varphi}{\partial t}
=-\frac{18(1+3t)}{r}
\]

\[
\frac{v^2}{2}
=
\frac{9(1+3t)^4}{2r^4}
\]

所以压力场为：

\[
\boxed{
p=p_\infty
+\frac{18\rho(1+3t)}{r}
-\frac{9\rho(1+3t)^4}{2r^4}
}
\]

\section*{七、三个定理的最短背法}

\subsection*{课件题图：书 4.12 与 Kelvin 环量}

\begin{center}
\includegraphics[width=0.93\linewidth]{ch5_images/slide_sphere_kelvin.png}
\end{center}

\begin{center}
\begin{tabular}{>{\raggedright\arraybackslash}p{3cm} >{\raggedright\arraybackslash}p{4.2cm} >{\raggedright\arraybackslash}p{6cm}}
\toprule
\textbf{定理} & \textbf{条件} & \textbf{结论}\\
\midrule
开尔文定理 &
无粘、正压、质量力有势、连续流场 &
沿封闭流体线的速度环量不随时间变化：
\[
\frac{D\Gamma}{Dt}=0
\]
\\
拉格朗日定理 &
满足开尔文定理条件 &
若某时刻无旋，则过去和以后都无旋。
\\
海姆霍兹第一定理 &
满足开尔文定理条件 &
组成涡面的流体质点永远组成涡面。
\\
海姆霍兹第二定理 &
满足开尔文定理条件 &
涡管强度不随时间变化。
\\
\bottomrule
\end{tabular}
\end{center}

\begin{warnbox}{背诵重点}
\begin{itemize}
  \item 开尔文定理管的是环量守恒。
  \item 拉格朗日定理管的是无旋保持。
  \item 海姆霍兹定理管的是涡线、涡面、涡管强度保持。
  \item 这些定理的共同条件都离不开：无粘、正压、质量力有势、连续流场。
\end{itemize}
\end{warnbox}

\section*{八、今天 2.5 小时复习安排}

\begin{center}
\begin{tabular}{p{2.5cm} p{10cm}}
\toprule
\textbf{时间} & \textbf{任务}\\
\midrule
30 分钟 &
看何枫第5章课件，重点标出：理想流体条件、兰姆方程、定常/非定常伯努利、开尔文定理。
\\
60 分钟 &
做三道核心题：环量题、书 4.12、U 型管振荡。每题按“条件判断--方程--代入--结果”写。
\\
40 分钟 &
不看答案重写：书 4.12、非定常伯努利积分式、伯努利成立条件。
\\
20 分钟 &
背三个定理的条件和结论，并口头解释“为什么非定常题多一项 \(\partial\varphi/\partial t\)”。
\\
\bottomrule
\end{tabular}
\end{center}

\section*{九、最后自测}

如果下面问题能答出来，第5章今天就算过关。

\begin{enumerate}
  \item 理想流体在固壁上的边界条件是什么？
  \item 伯努利方程什么时候沿流线成立？什么时候全流场成立？
  \item 非定常无旋伯努利比定常伯努利多哪一项？
  \item 环量的定义是什么？
  \item 书 4.12 为什么先用 \(4\pi r^2v_r=4\pi R^2\dot R\)？
  \item 开尔文定理成立条件是什么？
  \item 拉格朗日定理说明了什么？
  \item 涡管强度是什么意思？
\end{enumerate}

\end{document}
